\section{Conclusion}

\label{sec:conclusion}
%----------------------------------------------------------------------

DeFi's transparency is both its strength and its fatal flaw. The same public mempool that enables permissionless participation also enables systematic value extraction. \MEV{} is not a bug in DeFi---it is a consequence of plaintext orderflow on a public ledger.

Confidential DeFi on \Zoo{} Network eliminates this problem at the cryptographic level. Orders are encrypted before they enter the mempool. Matching happens on ciphertexts. Only matched pairs are decrypted. The result is a trading system where the fastest extractor has no advantage because there is nothing to extract.

The overhead is real: 490$\times$ for a single encrypted match compared to plaintext. But this comparison misses the point. The relevant comparison is not against a platonic frictionless market; it is against the actual DeFi market where users lose \$1.5 billion annually to \MEV{}. Against that baseline, a 2.8-second settlement with zero \MEV{} is an improvement.

Three open problems remain:

\begin{enumerate}
    \item \textbf{Cross-Chain Confidential DeFi.} Extending encrypted orderflow to atomic swaps across chains, where each chain holds a different \FHE{} key.
    \item \textbf{Encrypted Derivatives.} Applying \FHE{} to options and futures, where payoff calculations are more complex than spot matching.
    \item \textbf{Recursive FHE.} Using recursive bootstrapping to support unbounded-depth circuits, enabling complex matching algorithms (e.g., pro-rata allocation) without noise budget exhaustion.
\end{enumerate}

Front-running dies when there is nothing to see.

\begin{thebibliography}{20}

\bibitem{flashbots2024mev}
Flashbots,
``MEV-Explore: Quantifying Maximal Extractable Value,''
\url{https://explore.flashbots.net}, 2024.

\bibitem{chillotti2020tfhe}
I.~Chillotti, N.~Gama, M.~Georgieva, and M.~Izabach\`{e}ne,
``TFHE: Fast Fully Homomorphic Encryption over the Torus,''
\textit{J.\ Cryptology}, vol.~33, pp.~34--91, 2020.

\bibitem{cheon2017ckks}
J.~H.~Cheon, A.~Kim, M.~Kim, and Y.~Song,
``Homomorphic Encryption for Arithmetic of Approximate Numbers,''
\textit{Proc.\ ASIACRYPT}, Springer, 2017.

\bibitem{daian2020flash}
P.~Daian, S.~Goldfeder, T.~Kell, Y.~Li, X.~Zhao, I.~Bentov, L.~Breidenbach, and A.~Juels,
``Flash Boys 2.0: Frontrunning in Decentralized Exchanges, Miner Extractable Value, and Consensus Instability,''
\textit{Proc.\ IEEE S\&P}, 2020.

\bibitem{zoo-dex}
Z.~Kelling,
``Zoo DEX: Decentralized Exchange on Zoo Network,''
Zoo Labs Foundation, 2026.

\bibitem{zoo-fhe}
Z.~Kelling,
``Threshold Fully Homomorphic Encryption on Zoo Network: GPU-Accelerated Confidential Computation,''
Zoo Labs Foundation, 2026.

\bibitem{zoo-threshold-signatures}
Z.~Kelling,
``Threshold Signatures on Zoo T-Chain,''
Zoo Labs Foundation, 2026.

\bibitem{penumbra2023}
Penumbra Labs,
``Penumbra Protocol Specification,''
\url{https://protocol.penumbra.zone}, 2023.

\bibitem{renegade2023}
Renegade,
``Renegade: On-Chain Dark Pool Protocol,''
\url{https://renegade.fi/whitepaper.pdf}, 2023.

\bibitem{vickrey1961}
W.~Vickrey,
``Counterspeculation, Auctions, and Competitive Sealed Tenders,''
\textit{J.\ Finance}, vol.~16, no.~1, pp.~8--37, 1961.

\end{thebibliography}

